% Corrigé intégré automatiquement pour la banque Exercices.
% Source : A_Integrer/DDS_05/088_Geometrie/088_Geometrie.tex
\begin{corrigebox}[Corrigé]
\CorrigeQuestion{1}
On paramètre les mouvements relatifs, puis on applique la composition des vitesses et, si nécessaire, la dérivation vectorielle dans le repère adapté. Pour un train épicycloïdal, on écrit la relation de Willis avec les nombres de dents avant d'imposer les éléments bloqués ou entraînés.
\CorrigeQuestion{2}
On identifie les classes d'équivalence cinématique, les surfaces de contact et les mouvements relatifs autorisés. Chaque contact est remplacé par la liaison normalisée correspondante, puis les axes et points caractéristiques sont reportés sur le schéma cinématique minimal.
\CorrigeQuestion{3}
La résolution consiste à identifier les données et l'inconnue, choisir la loi du cours adaptée, écrire l'équation symbolique avant toute application numérique, puis vérifier l'homogénéité, le signe et l'ordre de grandeur du résultat. Les hypothèses utilisées doivent être explicitement contrôlées à la fin.
\CorrigeQuestion{4}
La résolution consiste à identifier les données et l'inconnue, choisir la loi du cours adaptée, écrire l'équation symbolique avant toute application numérique, puis vérifier l'homogénéité, le signe et l'ordre de grandeur du résultat. Les hypothèses utilisées doivent être explicitement contrôlées à la fin.
\CorrigeQuestion{5}
La résolution consiste à identifier les données et l'inconnue, choisir la loi du cours adaptée, écrire l'équation symbolique avant toute application numérique, puis vérifier l'homogénéité, le signe et l'ordre de grandeur du résultat. Les hypothèses utilisées doivent être explicitement contrôlées à la fin.
\CorrigeQuestion{6}
La résolution consiste à identifier les données et l'inconnue, choisir la loi du cours adaptée, écrire l'équation symbolique avant toute application numérique, puis vérifier l'homogénéité, le signe et l'ordre de grandeur du résultat. Les hypothèses utilisées doivent être explicitement contrôlées à la fin.
\CorrigeQuestion{7}
La résolution consiste à identifier les données et l'inconnue, choisir la loi du cours adaptée, écrire l'équation symbolique avant toute application numérique, puis vérifier l'homogénéité, le signe et l'ordre de grandeur du résultat. Les hypothèses utilisées doivent être explicitement contrôlées à la fin.
\CorrigeQuestion{8}
La résolution consiste à identifier les données et l'inconnue, choisir la loi du cours adaptée, écrire l'équation symbolique avant toute application numérique, puis vérifier l'homogénéité, le signe et l'ordre de grandeur du résultat. Les hypothèses utilisées doivent être explicitement contrôlées à la fin.
\end{corrigebox}
